\chapter{Outlook}
\label{Chpt:Outlook}


%The work presented in this thesis represents, to the best of my knowledge,  the most complete characterization of the interaction properties and Feshbach molecules of a mass-imbalanced Fermi-Fermi mixture of a lanthanide and an alkali to date. 

In this thesis, we have presented a thorough characterization of the interaction properties of a mass-imbalanced Fermi-Fermi mixture of a lanthanide and an alkali. Looking for a suitable, broad and well isolated interspecies Feshbach resonance, we identified more than $40$ resonances. Among these, we focused on the low magnetic field region below \SI{10}{G}, where we identified a Feshbach resonance close to \SI{7}{G}, which is the most promising candidate to realize and control strong interaction in the Dy-K mixture. This results are presented in Chapter~\ref{Chpt:Lowfieldpaper}. 

we have characterized 
the \SI{7}{G} resonance in detail, studied the strongly interacting regime 
through the collisional hydrodynamic crossover, in Chapter~\ref{Chpt:DipoleandHydro}, and produced and characterized optically trapped DyK Feshbach molecules, in Chapter~\ref{Chpt:MolTrap}, identifying and partially mitigating the loss processes that limit their lifetime, in Chapter~\ref{Chpt:MolLoss}. In Chapter~\ref{Chpt:Summary}, we presented an honest assessment of what we have learned, from a more technical point of view, in the perspective of a second generation experiment.  Here, I would like to look forward to the next possibilities our current experimental set-up offers.
%Here, we give a forward look, 



The most natural next step is the realization of a molecular Bose-Einstein 
condensate of DyK dimers. This would represent the bosonic limit of the 
\ac{BEC}-\ac{BCS} crossover in a mass-imbalanced system, and would open the door to the study of the full crossover from a \ac{mBEC} to the fermionic superfluid regime.  The results of Chapter~\ref{Chpt:MolLoss} represent significant progress toward this goal: by identifying trapping conditions where light-induced losses are substantially suppressed, we have extended the molecular lifetime to the point where further cooling becomes conceivable. Reaching a \ac{mBEC} will likely require a combination of improved trapping conditions 
 
exploiting longer wavelength light sources, as suggested by our results 
 
and a more degenerate starting mixture, which in turn requires addressing the evaporative cooling limitations discussed in Chapter~\ref{Chpt:Summary}. Once a sufficiently degenerate and long-lived molecular sample is available, the study of the full \ac{BEC}-\ac{BCS} crossover in a mass-imbalanced system becomes accessible, including the search for the exotic pairing phases, such as the \ac{FFLO} state, that motivated the construction of the experiment in the first place.

Beyond the molecular side of the crossover, the atomic mixture itself offers 
interesting directions that I believe are worth pursuing. The study of 
collective oscillations, and in particular the dipole-mode spectrum presented 
in Chapter~\ref{Chpt:DipoleandHydro}, was performed at a fixed degeneracy 
level set by the current experimental constraints. A natural extension would 
be to map out the temperature dependence of the hydrodynamic crossover, 
tracking how the onset of collisional hydrodynamics evolves as the mixture 
is brought closer to degeneracy. This would provide a more complete picture 
of the strongly interacting regime and could reveal signatures of many-body 
pairing around the superfluid transition temperature


, in a regime analogous 
to the pseudogap phase discussed in the context of conventional 
superconductors.





In the last few years, thanks to an increasing interest in exotic mixtures of alkali and lanthanides, some first theoretical results have started to become available, like for example the estimation of Van der Waals coefficients for Er or Dy mixed with alkali or alkali-earth \cite{Kopczyk2024vdw}. This is concurrent with more experiments developing in the last few years, mixtures of alkali with open-shell lanthanides are becoming more popular. The recent results in the field consist of observation of Feshbach resonances in various mixtures of Er-Li \cite{Schafer2022fro}, the preparation of doubly degenerate Fermi-Fermi \cite{Karpov2026ico} and Bose-Bose \cite{Kalia2026coa} mixtures of Er-Li, the first studies of Er-K (L.~Klaus, private communication), and the observation of Feshbach resonances in a Bose-Fermi Dy-Li mixture \cite{Xie2025fso}
 






